% RESEARCH_DRAFT_SHA256: a148d5285ca952d9b01a1914d671239de2b1eceba4a16e04b37156ebd421f6cc
% Layout and source-reference wrappers expanded for the web renderer.
\documentclass[11pt]{article}
\usepackage[margin=25mm]{geometry}
\usepackage{amsthm}
\usepackage{fontspec,unicode-math}
\setmainfont{Latin Modern Roman}
\setmathfont{Latin Modern Math}
\usepackage{amsmath,amssymb,mathtools,xurl,hyperref}
\hypersetup{colorlinks=true,urlcolor=blue,linkcolor=blue}
\setcounter{secnumdepth}{0}
\setlength{\parindent}{0pt}
\setlength{\parskip}{5pt}
\setlength{\emergencystretch}{3em}
\newcommand{\sref}[2]{\hyperlink{source:#1:#2}{[S#2]}}
\newcommand{\cataloguescope}[1]{\paragraph{Source record.}#1}
\newcommand{\floor}[1]{\left\lfloor#1\right\rfloor}
\newtheorem{proposition}{Proposition}
\newtheorem{lemma}[proposition]{Lemma}
\begin{document}
\paragraph{Research manuscript.} OpenAI \texttt{gpt-6-astra}, reasoning \texttt{ultra}; 11 October 2026.
\paragraph{Preserved source body.} The complete catalogue mathematical text and sources follow with layout and source links expanded, including source qualifications and its historical preparation notes. New research begins at the section headed Research attempt.

\section{Resolved Triple-Product Group Challenge}\label{statement}

\subsection{Definitions and mathematical statement}

\paragraph{Shared notation.}Standard number systems are $\mathbb N,\mathbb Z,\mathbb Q,\mathbb R,\mathbb C$. Specialized notation and index ranges follow the cited source and the definitions below.


For subsets $S,T,U$ of a finite group, the triple-product property requires $s_1s_2^{-1}t_1t_2^{-1}u_1u_2^{-1}=1$ only when each selected pair is equal. The matrix-multiplication exponent is the infimum of exponents achievable by the stated arithmetic model.
\paragraph{Record scope and qualifications.}
The catalogue title is retained as metadata. Named certificates, contracts, bounds or frameworks in that title are not asserted as results of the original source extract.
The supplied registry flags this extraction as unrecoverable. It is displayed as a literal source fragment; the definitions provide context rather than a replacement formal statement.

The registry flags the extraction as unrecoverable and its title calls the challenge resolved; the subset sizes, representation-degree conditions, and target exponent are not supplied.

Source-record limitation: Source flags the statement as unrecoverable: The official report poses this challenge and then constructs such groups in the same abstract, obtaining nontrivial bounds on the matrix-multiplication exponent. No separate clean open statement is promoted; the original remains retained for provenance.

The fragment does not specify the conditions needed to turn the named challenge into an exact remaining exponent question.
The source text below is rendered literally where necessary; damaged formulas and composite extracts are not silently repaired into a new mathematical claim.
The supplied catalogue marks this record as unrecoverable. Its extract and terminology are documented, but a complete original problem has not been reconstructed.
\paragraph{Statement recorded in the supplied source.}
Find families of groups with the triple product property and character degrees sufficiently small to yield nontrivial bounds on the matrix-multiplication exponent. \href{https://doi.org/10.4171/owr/2005/26}{[S2]}
\paragraph{Source status.}
The supplied catalogue labels this source record open; that metadata was checked against the supplied snapshot on 2026-10-11. The label is preserved for provenance and is not an independent verification of present open status or mathematical correctness. The catalogue literature assessment is dated 2026-08-21.
\subsection{Short English statement}
Seek group-based constructions that improve matrix-multiplication complexity.
\subsection{Sources}
\begin{itemize}
\item[S1] Ulam AI and the UnsolvedMath Contributors, UnsolvedMath: A Curated Collection of Open Mathematics Problems, catalogue entry. Catalogue curation and metadata: CC BY 4.0. The source extract is retained; definitions, English summaries and qualifications are editorial additions.
\url{https://huggingface.co/datasets/ulamai/UnsolvedMath}.
\item[S2] Complexity Theory (2006), Oberwolfach Reports, DOI 10.4171/owr/2005/26, pp. 1431--1504
\url{https://doi.org/10.4171/owr/2005/26}.
\item[S3] Cohn--Umans group approach
\url{https://arxiv.org/abs/math/0307210}.
\item Source-specific Oberwolfach Reports licences, including CC BY-SA 4.0 where applicable, are retained; see THIRD\_PARTY\_NOTICES.md.
\end{itemize}
\paragraph{Preparation and review.}
The original source and UnsolvedMath contributors retain attribution for the source extract. Codex (GPT-6; exact serving revision not exposed) assisted with definitions, English summaries, scope notes and formatting. No new human mathematical authorship or independently verified solution is asserted. Catalogue identity checked against the complete supplied registry and current repository. Distinct upstream records, including repeated questions and broad or damaged extracts, are retained for complete coverage of the source CS category. A repeated record is not a new independent conjecture. Complete mathematical-equivalence and current-literature review has not been performed. Preparation and final review follow contribution prompts C and D, with source gaps disclosed. Structural and rendering checks do not establish a new theorem or open status.
\label{end:statement}

\section{Research attempt}
\subsection{Literature, source context and scope}
The full original catalogue body and source qualifications are preserved above. The relevant primary workshop passages were read in the cached publisher report. Historical claims and subsequent attributed results are distinguished from the mathematical work below. A restricted example, standard lemma or known theorem is not counted as new progress on the entire original target. Missing assumptions and repeated or rhetorical extracts are not silently turned into different solved problems.
Codex (GPT-6; exact serving revision not exposed) supplied the auxiliary proof route used here. OpenAI \texttt{gpt-6-astra}, reasoning \texttt{ultra}, independently worked through its mathematical steps, wrote this exposition, and performed the finite checks. This is model collaboration, not independent human verification.

\subsection{An abelian-group obstruction to one proposed route}
The source fragment is explicitly unrecoverable as a clean remaining problem: the same primary abstract presents nonabelian constructions and exponent improvements. We do not invent a new exponent target. Instead, test whether restricting to abelian groups, whose representations are small, can by itself yield a better-than-cubic bound through one triple.
\begin{proposition}
If $G$ is a finite abelian group and $S,T,U\subseteq G$ have the stated triple-product property, then $|S||T||U|\le|G|$.
\end{proposition}
\begin{proof}
Consider $(s,t,u)\mapsto stu$. If $s_1t_1u_1=s_2t_2u_2$, commutativity permits rearrangement to $s_1s_2^{-1}t_1t_2^{-1}u_1u_2^{-1}=1$. The triple-product property forces equality in each coordinate, so the map is injective. Counting its domain and codomain proves the inequality.
\end{proof}
The estimate is tight: in $G=(\mathbb Z/k\mathbb Z)^3$, take the three coordinate axes as $S,T,U$. The quotient-product equation becomes three separate coordinate equalities, so the property holds, and all sets have size $k$, giving $|G|=k^3$.
In a finite-dimensional complex irreducible representation of an abelian group, all representing matrices commute. A commuting finite family has a common eigenvector: take an eigenspace of the first matrix, invariant under the others, and successively restrict to an eigenspace of each remaining matrix within the invariant space. The resulting nonzero common eigenspace contains a common eigenvector. Its line is invariant under the group, so irreducibility forces dimension one. Using the primary report's representation identity $\sum_i d_i^2=|G|$ and its attributed group-algebra inequality $k^\omega\le\sum_i d_i^\omega$, the right side for an abelian group is therefore $|G|$. For equal set sizes $k>1$, the bound on $\omega$ obtained from this single inequality is $\log_k|G|\ge3$ by the proposition, so this route alone supplies no subcubic exponent bound.
\subsection{Verification and remaining gaps}
The coordinate-axis triples and injectivity were exhaustively checked for $k=2,3$. The representation identity and group-algebra complexity inequality are source results used explicitly as context, not proved anew or advertised as our work. The attempted inference that small character degrees alone suffice fails because the group must also accommodate the triple. This does not rule out nonabelian groups, simultaneous triples, different embeddings or the constructions already in the source. No new matrix-multiplication exponent or repaired remaining challenge is claimed.

\section{References and source attribution}
Primary workshop report: \url{https://doi.org/10.4171/owr/2005/26}. Relevant passages read: 2005/26 pp. 1469--1471, triple-product embeddings, character-degree inequality and constructions already given.

All retrievals were checked on 11 October 2026. Original question attribution remains in the preserved Sources section. External theorems are attributed and not claimed as new discoveries or fully verified merely by citation.


\paragraph{Integration review.} Codex (GPT-6; exact serving revision not exposed) reviewed the written arguments and expanded layout and source-reference wrappers. Repository metadata and the production web renderer were checked using Python 3.13 and JavaScript with MathJax 3.2.2. This model-assisted review is not independent human proof verification; no PDF compilation or formal proof checking is asserted.

\section{Completion Estimate}
\noindent\textbf{COMPLETION: 0\%}
No substantiated new progress toward the entire original target is claimed. The worked results and failed route are delimited above; a known result, illustrative model, or auxiliary obstruction is not counted as original-target completion.

\section{Final status and attribution}
\textbf{Attempt status: unresolved.} This describes this research attempt, not a new ruling on the source catalogue's status or on external literature claims.
The mathematics and this self-audit were generated by OpenAI \texttt{gpt-6-astra}, reasoning \texttt{ultra}, on 11 October 2026. Codex (GPT-6; exact serving revision not exposed) orchestrated the batch and contributed formatting/repository review. No human mathematical authorship, independent human proof checking, or formal verification is claimed. The historical preparation paragraphs in the reproduced catalogue body concern that body only.

Contribution prompt A was applied to the complete definition and supplied source record, followed by the adversarial checks of prompt B. Local reproduction files are \texttt{batch016/checks.py}, \texttt{computations.json}, and the per-ID audit JSON. They were not public links at preparation time. Numerical checks are finite and do not prove unrestricted results. No LaTeX engine was found locally; compilation was not performed by this research agent.

\end{document}
